\section{Related Work}

\subsection{Spurious Correlation Benchmarks}

\textbf{Waterbirds} \cite{sagawa2019distributionally} correlates bird species with background (waterbirds appear on water 90\% of training samples). Standard ERM training achieves 95\% average accuracy but 60\% worst-group accuracy, demonstrating shortcut learning. \textbf{CelebA} \cite{liu2015faceattributes} exhibits gender-hair color correlation (blonde hair occurs predominantly in female training images). \textbf{Spawrious} \cite{lynch2023spawrious} provides synthetic benchmarks with controllable correlation rates (50\%-95\%).

These benchmarks demonstrate that spurious features harm minority group performance, but existing detection methods require knowing which feature is spurious. We address unknown spurious detection via gradient abnormality.

\subsection{Robust Learning Methods}

\textbf{GroupDRO} \cite{sagawa2019distributionally} minimizes worst-group loss via group annotations, achieving 80-85\% WGA on Waterbirds (294 GitHub stars, widely reproduced). Requires explicit group membership during training.

\textbf{JTT} \cite{liu2021just} trains a two-stage model: (1) identify hard samples from initial ERM, (2) upsample and retrain. Achieves $\sim$78\% WGA on Waterbirds (72 GitHub stars). Requires no annotations but assumes initial model correctly identifies minority samples.

\textbf{SCER} \cite{park2025spurious} refines prototypes in embedding space to correct spurious predictions, reporting $\sim$90\% WGA (estimated, code unavailable for reproduction). Operates post-hoc and misses gradient-level training dynamics.

\textbf{Gap:} GroupDRO/JTT require annotations or two-stage training. SCER operates on embeddings, not gradients. We propose gradient-based detection + mitigation in a unified single-stage framework.

Table~\ref{tab:comparison} summarizes key differences:

\begin{table}[t]
\centering
\caption{Comparison of spurious correlation mitigation methods}
\label{tab:comparison}
\begin{tabular}{lcccc}
\toprule
Method & WGA & Annotations? & Stages & Space \\
\midrule
GroupDRO \cite{sagawa2019distributionally} & 80-85\%* & Yes & Single & Loss \\
JTT \cite{liu2021just} & $\sim$78\%* & No & Two & Sample \\
SCER \cite{park2025spurious} & $\sim$90\%* & No & Post-hoc & Embedding \\
\textbf{Ours (PoC)} & 78\%** & No & Single & Gradient \\
\bottomrule
\end{tabular}
\\[0.5em]
\small{*Literature-reported. **MNIST smoke test (1 seed, 2 epochs), not Waterbirds.}
\end{table}

\subsection{Gradient Attribution for Spurious Detection}

\textbf{Grad-CAM} \cite{selvaraju2017grad} produces saliency maps via weighted activation gradients. Widely used for interpretability (10,000+ citations) but fails for unknown spurious: user must know where to look.

\textbf{Integrated Gradients} \cite{sundararajan2017axiomatic} computes attribution via path integral from baseline to input. Adebayo et al.~(2022) \cite{adebayo2022post} showed via user study (109 cites) that practitioners cannot detect unknown spurious correlations using IG or Grad-CAM results. \textbf{Key insight:} Attribution explains \emph{which} features matter (result interpretation), not \emph{whether} the process is abnormal (mechanism detection).

\textbf{GAIA} (Gradient Abnormality for OOD) \cite{chen2023exploring} detects distribution shift via gradient zero-deflation (GAIA-Z) and channel variance (GAIA-A). Achieves FPR95 reduction of 45.41\% on CIFAR100 for OOD detection (16 cites). Operates on gradient \emph{process}, not \emph{result}.

\textbf{SPROD} \cite{lee2025detecting} detects spurious-correlated OOD via prototype-based divergence (4 cites). Complements our approach (prototypes vs gradients).

\textbf{Gap:} GAIA applied to OOD only, not subpopulation shift. SPROD post-hoc detection only. Adebayo shows attribution fails for unknown spurious. We extend GAIA to minority group detection (subpopulation shift within ID distribution) and propose gradient regularization for mitigation.

\subsection{Positioning Summary}

Existing work addresses spurious correlations via (1) annotation-based reweighting (GroupDRO), (2) two-stage training (JTT), or (3) post-hoc embedding refinement (SCER). Attribution methods (Grad-CAM, IG) fail for unknown spurious \cite{adebayo2022post}. GAIA detects OOD via gradient abnormality but not minority groups.

\textbf{Our contribution:} Extend GAIA to subpopulation shift, validate spurious-conflict mechanism via synthetic causality tests, and propose spatial gradient regularization as annotation-free mitigation. Complements embedding-based methods (SCER) by operating on gradient space.
